\section{Objectifs}
\begin{itemize}
\item Mesurer un angle en radians et utiliser la longueur d'arc.
\item Comprendre l'enroulement de la droite réelle sur le cercle trigonométrique.
\item Lire sinus et cosinus d'un réel et retrouver les valeurs remarquables.
\end{itemize}

\section{Cours}
\subsection{Radian et cercle trigonométrique}
Le cercle trigonométrique est le cercle de centre O et de rayon 1 orienté dans le sens direct. Un angle de \texttt{1} radian intercepte sur un cercle de rayon 1 un arc de longueur 1. Un tour complet mesure \texttt{2π} radians.

\subsection{Enroulement de la droite}
À chaque réel \texttt{x}, on associe un point du cercle en parcourant une longueur algébrique \texttt{x} depuis le point \texttt{(1,0)}. Des réels qui diffèrent d'un multiple de \texttt{2π} ont la même image.

\subsection{Sinus et cosinus}
Si le point image de \texttt{x} a pour coordonnées \texttt{(cos x, sin x)}, alors \texttt{cos²x+sin²x=1}. Les définitions prolongent celles du triangle rectangle.

Valeurs remarquables à connaître et retrouver par lecture du cercle : \texttt{π/6}, \texttt{π/4}, \texttt{π/3}, \texttt{π/2}, puis leurs angles associés. Par exemple \texttt{cos(π/4)=sin(π/4)=√2/2}, \texttt{cos(π/3)=1/2} et \texttt{sin(π/3)=√3/2}.

\subsection{Angles associés}
Les symétries du cercle donnent par exemple \texttt{cos(-x)=cos x}, \texttt{sin(-x)=-sin x}, \texttt{cos(π-x)=-cos x}, \texttt{sin(π-x)=sin x}.

\section{Démonstration à travailler}
Le programme demande notamment de savoir établir les valeurs de \texttt{cos(π/4)}, \texttt{sin(π/4)}, \texttt{cos(π/3)} et \texttt{sin(π/3)}.

\section{Algorithmique}
L'approximation de \texttt{π} par la méthode d'Archimède peut servir d'activité algorithmique ; elle sera surtout reliée au bloc algorithmique transversal en C22.

\section{Erreurs fréquentes}
\begin{itemize}
\item Confondre degrés et radians.
\item Inverser les coordonnées : sur le cercle le point est \texttt{(cos x, sin x)}.
\item Oublier la périodicité de l'enroulement.
\end{itemize}

\section{Synthèse}
Le cercle trigonométrique transforme tout réel en un point du cercle unité ; ses coordonnées définissent cosinus et sinus et permettent de retrouver les valeurs et signes des angles associés.
